% BleModule — BLE通信モジュール
\subsection{BleModule}

\begin{apibox}{BleModule --- BLE通信モジュール}
\textbf{定義ファイル:} \texttt{lib/BleModule/BleModule.h / .cpp}\\
\textbf{分類:} 入出力モジュール（\texttt{updateInput()} + \texttt{updateOutput()}）\\
\textbf{対象デバイス:} ESP32 BLE（GATT Server）

BLE GATT Serverとして動作し、RX/TX Characteristicを使用してデータの受信・送信を行う。
BLEコールバックは別タスクで実行されるため、volatileフラグ方式でスレッドセーフにデータを受け渡す。
\end{apibox}

\begin{cautionbox}[スレッドセーフティ]
BLEのコールバック（\texttt{onConnect}、\texttt{onDisconnect}、\texttt{onReceive}）はFreeRTOSの別タスクで実行される。
コールバックから\texttt{SystemData}に直接書き込むことは禁止。
内部バッファ + \texttt{volatile}フラグ方式で\texttt{updateInput()}を経由してデータを受け渡す。
\end{cautionbox}

\subsubsection{機能概要}

\begin{itemize}
    \item BLE GATT Serverの初期化・アドバタイズ開始
    \item RX Characteristic経由のデータ受信（コールバック→内部バッファ→SystemData）
    \item TX Characteristic経由のデータ送信（SystemData→notify）
    \item 接続状態の監視と切断時の自動リアドバタイズ
    \item \texttt{deinit()}によるBLEリソース解放
\end{itemize}

% --- 定数 ---
\begin{funcblock}{定数}
\begin{tabularx}{\textwidth}{l l X}
\toprule
\textbf{名前} & \textbf{値} & \textbf{説明} \\
\midrule
\texttt{BLE\_RX\_BUFFER\_SIZE} & \texttt{128} & 受信バッファサイズ [bytes] \\
\texttt{BLE\_TX\_BUFFER\_SIZE} & \texttt{128} & 送信バッファサイズ [bytes] \\
\bottomrule
\end{tabularx}
\end{funcblock}

% --- Config ---
\subsubsection{Config構造体}

\begin{funcblock}{BleConfig}
BLEサービスの識別情報を定義する。

\begin{longtable}{l l l p{4.5cm}}
\toprule
\textbf{メンバ} & \textbf{型} & \textbf{制約・既定値} & \textbf{説明} \\
\midrule
\endhead
\texttt{deviceName}  & \texttt{const char*} & ---  & BLEデバイス名（アドバタイズ時に表示） \\
\texttt{serviceUuid} & \texttt{const char*} & ---  & GATTサービスUUID文字列 \\
\texttt{rxCharUuid}  & \texttt{const char*} & ---  & RX Characteristic UUID（クライアント→デバイス方向） \\
\texttt{txCharUuid}  & \texttt{const char*} & ---  & TX Characteristic UUID（デバイス→クライアント方向） \\
\bottomrule
\end{longtable}
\end{funcblock}

% --- Data ---
\subsubsection{Data構造体}

\begin{funcblock}{BleData}
BLE通信の状態と送受信データを保持する。

\begin{longtable}{l l l p{3.8cm}}
\toprule
\textbf{メンバ} & \textbf{型} & \textbf{初期値} & \textbf{説明} \\
\midrule
\endfirsthead
\toprule
\textbf{メンバ} & \textbf{型} & \textbf{初期値} & \textbf{説明} \\
\midrule
\endhead
\midrule
\multicolumn{4}{r}{\small 次ページに続く} \\
\endfoot
\bottomrule
\endlastfoot
\multicolumn{4}{l}{\textbf{--- 接続状態 ---}} \\
\texttt{connected}   & \texttt{bool}    & \texttt{false} & BLE接続状態 \\
\midrule
\multicolumn{4}{l}{\textbf{--- 受信（入力フェーズで更新） ---}} \\
\texttt{newDataIn}   & \texttt{bool}    & \texttt{false} & 新規受信データあり。ロジックフェーズで参照後クリアする \\
\texttt{rxData[128]} & \texttt{uint8\_t[]} & \texttt{\{\}} & 受信データバッファ \\
\texttt{rxLength}    & \texttt{uint8\_t} & \texttt{0}    & 受信データ長 [bytes] \\
\midrule
\multicolumn{4}{l}{\textbf{--- 送信（ロジックフェーズで設定） ---}} \\
\texttt{sendRequest} & \texttt{bool}    & \texttt{false} & 送信リクエスト。\texttt{true}で\texttt{updateOutput()}がデータを送信 \\
\texttt{txData[128]} & \texttt{uint8\_t[]} & \texttt{\{\}} & 送信データバッファ \\
\texttt{txLength}    & \texttt{uint8\_t} & \texttt{0}    & 送信データ長 [bytes] \\
\end{longtable}
\end{funcblock}

% --- メソッド ---
\subsubsection{メソッド}

\begin{funcblock}{BleModule(const BleConfig\& config)}
\begin{tabularx}{\textwidth}{l X}
\toprule
\textbf{項目} & \textbf{内容} \\
\midrule
引数   & \texttt{config} --- \texttt{BleConfig}へのconst参照 \\
動作   & Configを\texttt{\_config}にコピー保持する \\
\bottomrule
\end{tabularx}
\end{funcblock}

\begin{funcblock}{bool init()}
\begin{tabularx}{\textwidth}{l X}
\toprule
\textbf{項目} & \textbf{内容} \\
\midrule
戻り値     & \texttt{bool} --- 常に\texttt{true} \\
動作       & BLEデバイス初期化、GATTサーバー作成、サービス・Characteristic登録、コールバック設定、アドバタイズ開始 \\
\bottomrule
\end{tabularx}
\end{funcblock}

\begin{funcblock}{void updateInput(SystemData\& data)}
\begin{tabularx}{\textwidth}{l X}
\toprule
\textbf{項目} & \textbf{内容} \\
\midrule
書き込み先     & \texttt{data.ble} \\
タイミング制御 & 毎ループ実行 \\
動作           & \texttt{volatile}フラグ\texttt{\_dataReceived}を確認。\texttt{true}なら内部バッファから\texttt{data.ble.rxData}にコピーし、\texttt{newDataIn = true}に設定。接続状態を\texttt{data.ble.connected}に反映 \\
\bottomrule
\end{tabularx}
\end{funcblock}

\begin{funcblock}{void updateOutput(SystemData\& data)}
\begin{tabularx}{\textwidth}{l X}
\toprule
\textbf{項目} & \textbf{内容} \\
\midrule
読み取り元     & \texttt{data.ble} \\
タイミング制御 & 毎ループ実行 \\
動作           & \texttt{sendRequest}が\texttt{true}かつ接続中であれば、\texttt{txData}をTX Characteristicにnotifyで送信。\texttt{sendRequest}を\texttt{false}にクリア \\
\bottomrule
\end{tabularx}
\end{funcblock}

\begin{funcblock}{void deinit()}
\begin{tabularx}{\textwidth}{l X}
\toprule
\textbf{項目} & \textbf{内容} \\
\midrule
動作   & BLEアドバタイズを停止し、BLEデバイスを終了処理する \\
\bottomrule
\end{tabularx}
\end{funcblock}

\begin{funcblock}{void onConnect() / void onDisconnect() / void onReceive(const uint8\_t* payload, uint8\_t length)}
\begin{tabularx}{\textwidth}{l X}
\toprule
\textbf{項目} & \textbf{内容} \\
\midrule
呼び出し元 & BLEコールバック（別タスクで実行される） \\
動作       & \texttt{onConnect}: \texttt{\_clientConnected = true}。\texttt{onDisconnect}: \texttt{\_clientConnected = false}、リアドバタイズ開始。\texttt{onReceive}: 内部バッファにコピーし\texttt{\_dataReceived = true} \\
\bottomrule
\end{tabularx}
\end{funcblock}

\begin{notebox}[volatileフラグ方式のデータフロー]
BLEコールバック（別タスク）からSystemDataへのデータ受け渡しは以下の流れで行う:

\begin{enumerate}
    \item \textbf{コールバック}: 受信データを内部バッファ\texttt{\_rxBuffer}にコピーし、\texttt{volatile bool \_dataReceived = true}
    \item \textbf{updateInput()}: \texttt{\_dataReceived}を確認。\texttt{true}なら内部バッファから\texttt{data.ble.rxData}にコピーし、\texttt{\_dataReceived = false}にリセット
    \item \textbf{ロジックフェーズ}: \texttt{data.ble.newDataIn}を確認して受信データを処理
\end{enumerate}
\end{notebox}

% --- 使用例 ---
\subsubsection{使用例}

\begin{lstlisting}[style=cppstyle, caption={BleModuleの使用例}]
// ProjectConfig.h
const BleConfig BLE_CONFIG = {
    .deviceName  = "ESP32-Module",
    .serviceUuid = "12345678-1234-1234-1234-123456789abc",
    .rxCharUuid  = "12345678-1234-1234-1234-123456789abd",
    .txCharUuid  = "12345678-1234-1234-1234-123456789abe",
};

// ロジックフェーズでの使用
void applyPattern(SystemData& data) {
    // 受信データの処理
    if (data.ble.newDataIn) {
        // rxDataを解析
        data.ble.newDataIn = false;  // フラグクリア
    }
    // センサーデータのBLE送信
    if (data.ble.connected && data.mpu.isValid) {
        memcpy(data.ble.txData, &data.mpu.accelX,
               sizeof(float));
        data.ble.txLength    = sizeof(float);
        data.ble.sendRequest = true;
    }
}
\end{lstlisting}
